\documentclass[11pt]{article}

% --- Math and algorithms ---
\usepackage{amsmath, amssymb, amsthm, amsfonts}
\usepackage{mathrsfs}
\usepackage{algorithm}
\usepackage{algpseudocode}

% --- Encoding and layout ---
\usepackage[utf8]{inputenc}
\usepackage[T1]{fontenc}
\usepackage{setspace}
\usepackage{float}
\usepackage{graphicx}
\usepackage[table]{xcolor}
\usepackage{booktabs}
\usepackage{array}
\usepackage{tabularx}
\usepackage{multirow}
\usepackage{makecell}
\usepackage{ragged2e}
\usepackage{siunitx}
\usepackage[authoryear]{natbib}

\newcolumntype{Y}{>{\RaggedRight\arraybackslash}X}
\pagestyle{plain}
\setlength{\textwidth}{6.5in}
\setlength{\oddsidemargin}{0in}
\setlength{\evensidemargin}{0in}
\setlength{\textheight}{9.2in}
\setlength{\topmargin}{-0.2in}
\setlength{\headheight}{0in}
\setlength{\headsep}{0in}
\setlength{\footskip}{.4in}
\renewcommand{\baselinestretch}{1.25}

\newtheorem{theorem}{Theorem}
\newtheorem{obs}{Observation}
\newtheorem{proposition}{Proposition}
\newtheorem{remark}{Remark}
\newtheorem{definition}{Definition}

\begin{document}

\begin{center}
{\Large \bf Integrated Sample Size Optimization and Test Chamber Scheduling under Environmental Constraints}\\
\vskip 12pt
\normalsize
Nur Per\c{c}in$^{1}$,
Ali Ekici$^{1}$\\[6pt]
\footnotesize
$^{1}$Department of Industrial Engineering, Ozyegin University, Istanbul, Turkiye\\[3pt]
\texttt{nur.percin@ozu.edu.tr, ali.ekici@ozyegin.edu.tr}
\end{center}

\begin{abstract}
Research and development testing of refrigerators requires the coordinated use of heterogeneous test chambers under temperature, humidity, voltage, stage-precedence, and due-date constraints. In current practice, sample sizes are often fixed, although sample decisions directly determine both workload and opportunities for parallel execution. This study introduces the Test Chamber Scheduling with Sample Optimization Problem (TCSSOP), in which product sample sizes, chamber-environment assignments, and resource-feasible schedules are optimized jointly.

An iterative two-stage framework is proposed. For a given sample vector, chamber environments are assigned by balancing environment-specific and voltage-sensitive workload over compatible chamber capacity. Given chamber settings, sample sizes are improved by local search and evaluated through an earliest-due-date constructive scheduler with explicit feasibility validation. When the search stagnates, a Variable Neighborhood Search (VNS) mechanism perturbs sample sizes using fixed fallback ratios and re-optimizes the resulting workload structure.

Computational experiments on 90 benchmark scenarios show that endogenous sample optimization substantially reduces total tardiness relative to the fixed three-sample policy. The best configuration, using a \(0.20\)--\(0.30\) VNS fallback policy with \(K_{\max}=200\), provides the strongest overall performance. The results show that sample allocation should be treated as a tactical planning decision and coordinated with chamber reconfiguration rather than fixed uniformly across projects.
\end{abstract}

\textbf{Keywords:} Test chamber scheduling; Sample size optimization; Decision-dependent workload; Resource-constrained scheduling; Variable neighborhood search; Multi-stage testing systems

\section{Introduction}
\label{sec:intro}

Refrigerator manufacturers operate under compressed product-development cycles, strict market-specific compliance requirements, and high pressure to launch new models on time. Before commercialization, products must complete design approval tests covering performance, safety, energy efficiency, and consumer-usage conditions \citep{IEA2021,IEC60335-1,IEC60335-2-24,IEC62552}. Delays in these tests postpone certification and production release, while inefficient use of chambers lowers the return on capital-intensive R\&D infrastructure \citep{OECD_MSTI2024,OECD_MSTI2025}.

The testing system considered in this study is motivated by refrigerator R\&D operations at Vestel. Products compete for a set of heterogeneous climate chambers, each containing several stations that operate under a common environmental configuration. Test operations require specific temperature and humidity conditions; some products also require voltage-compatible execution. In addition, the testing process follows stage-state logic: gas amount determination precedes pull-down testing, downstream tests are released after the pull-down phase, and final consumer-usage tests require the relevant downstream stage conditions to be satisfied.

A key operational feature is the use of physical samples. Industrial practice commonly applies a fixed three-sample policy. This rule is simple, but it ignores the fact that sample size is itself a scheduling decision. Pull-down tests scale with the number of samples and therefore increase upstream workload, whereas additional samples can enable downstream tests to run in parallel and reduce completion times. The effect of sample size on tardiness is consequently non-monotone and depends on due dates, chamber eligibility, environmental bottlenecks, and stage-state interactions.

We define the resulting problem as the \emph{Test Chamber Scheduling with Sample Optimization Problem} (TCSSOP). TCSSOP combines heterogeneous parallel resources, chamber-level environmental assignment, humidity and voltage eligibility, multi-stage test logic, and decision-dependent workload generation. The objective is to minimize total tardiness, with total sample usage used as a secondary lexicographic criterion.

This study makes three contributions. First, it benchmarks the fixed three-sample policy against endogenous sample-size optimization. Second, it develops an iterative framework that coordinates sample-size decisions with chamber-environment reassignment and feasible scheduling. Third, it introduces a VNS diversification strategy that perturbs workload-defining sample decisions rather than only sequencing or assignment decisions. The computational results show that these components jointly produce substantial tardiness reductions across 90 benchmark scenarios.

\section{Literature Review}
\label{sec:literature_review}

The problem is related to several scheduling and test-planning streams. Resource-constrained multi-project scheduling and unrelated parallel-machine scheduling provide the closest scheduling background because multiple projects compete for heterogeneous resources under precedence, due-date, and eligibility restrictions. Recent surveys emphasize that practical scheduling systems increasingly require integrated treatment of resource heterogeneity, precedence, and operational constraints \citep{HartmannBriskorn2022,GomezSanchezLallaRuizFernandezGilCastroVoss2023}. In unrelated parallel-machine environments, machine eligibility and tardiness-oriented objectives substantially increase computational difficulty and motivate specialized heuristics \citep{UlagaDurasevicJakobovic2022,MaeckerShenMonch2023}.

Industrial scheduling often also involves setup or configuration decisions. The survey by \citet{Allahverdi2015} shows that ignoring setup-related effects can produce unrealistic schedules. In test-chamber systems, chamber environmental settings play a similar role because each chamber can serve only tests compatible with its current temperature-humidity configuration. This creates a reconfigurable resource-allocation problem related in spirit to reconfigurable manufacturing systems \citep{Koren1999,MehrabiUlsoyKoren2000}, although the present setting concerns testing rather than production.

Integrated decision-making models such as order acceptance and scheduling combine workload selection with sequencing under capacity limits \citep{Slotnick2011,WangYe2019}. However, the workload associated with an accepted job is usually exogenous. Reliability demonstration testing, in contrast, optimizes sample size and test duration under statistical risk or cost criteria \citep{Fernandez2022}, but typically abstracts away detailed chamber capacity and stage-dependent scheduling. TCSSOP lies between these areas: sample sizes generate the workload, while the resulting operations must be scheduled on constrained heterogeneous chambers.

Methodologically, VNS and related neighborhood-based metaheuristics have performed well in complex scheduling settings because they combine local improvement with systematic diversification \citep{MejiaYuraszeck2020}. The diversification mechanism proposed here differs from conventional VNS designs by perturbing sample-size decisions. Thus, the search explores alternative workload structures, which then induce different chamber-utilization and scheduling regimes.

\section{Problem Definition}
\label{sec:pd}

\subsection{System description}
\label{subsec:system}

The facility contains a set of chambers \(\mathcal{C}\). Chamber \(c\) has \(M_c\) stations and supports temperature adjustment; only selected chambers support humidity adjustment \(h_c\) and voltage adjustment \(v_c\). Stations in the same chamber operate under one common environment \(X_c\), where an environment represents a temperature-humidity combination.

The set of products is denoted by \(\mathcal{P}\). Each product \(p\) has a due date \(d_p\), a voltage requirement indicator \(\nu_p\), and a product-specific required test set \(\mathcal{T}_p \subseteq \mathcal{T}\). Test requirements are determined by target market, climatic class, energy-efficiency regulation, and electrical specification. Table~\ref{tab:test_attributes_short} summarizes the test catalog used in the study.

\begin{table}[H]
\caption{Test catalog and environmental requirements.}
\label{tab:test_attributes_short}
\centering
\footnotesize
\setlength{\tabcolsep}{3.2pt}
\begin{tabular}{|c|c|l|c|c|c|c|}
\hline
\textbf{Stage} & \textbf{No.} & \textbf{Test name} & \textbf{Temp.} & \textbf{Humidity} & \textbf{Samples} & \textbf{Days} \\ \hline
1 & 1 & Gas Amount Determination & 43/38/32~\si{\degreeCelsius} & Ambient & 1 & 10 \\ \hline
2 & 2 & Pull-down & 43/38/32~\si{\degreeCelsius} & Ambient & All & 5 \\ \hline
3 & 3--5 & Energy Efficiency & 16/25/32~\si{\degreeCelsius} & Ambient & 1 & 15 \\ \hline
3 & 6--11 & Performance & 10/16/25/32/38/43~\si{\degreeCelsius} & Ambient & 1 & 10 \\ \hline
4 & 12--13 & Freezing Capacity / Temperature Rise & 25~\si{\degreeCelsius} & Ambient & 1 & 10/5 \\ \hline
5 & 14--16 & Consumer Usage & 10/25/32~\si{\degreeCelsius} & Ambient or 85\% RH & 1 & 15 \\ \hline
\end{tabular}
\vspace{1mm}

\begin{minipage}{0.92\textwidth}
\footnotesize
Some energy and performance variants are alternatives selected according to climatic class and market regulation. Freezing capacity and temperature rise are optional.
\end{minipage}
\end{table}

The chamber system includes 28 chambers and 240 stations. Chamber capacities range from 6 to 12 stations. Humidity and voltage capabilities are heterogeneous, so the effective capacity available for a given environment may be much smaller than the total station count. Figure~\ref{fig:chamber_3d_short} illustrates the chamber-station structure.

\begin{figure}[H]
\centering
\includegraphics[width=0.88\textwidth]{chamber_structure.png}
\caption{Test chambers consist of multiple stations operating under a shared chamber-level environment.}
\label{fig:chamber_3d_short}
\end{figure}

\subsection{Decision-dependent workload}
\label{subsec:workload}

Each product has a sample-size decision \(s_p\), bounded by \(s_{\min} \le s_p \le s_{\max}\). Pull-down tests are executed once per sample, whereas other required tests are executed once per product. Thus, the multiplicity of test \(t\) for product \(p\) is
\[
m_{pt}(s_p)=
\begin{cases}
s_p, & \text{if }t\text{ is a pull-down test,}\\
1, & \text{otherwise.}
\end{cases}
\]
For environment \(e\), the induced workload is
\[
W_e(S)=
\sum_{p\in\mathcal{P}}
\sum_{\substack{t\in\mathcal{T}_p\\e(t)=e}}
m_{pt}(s_p)\tau_t,
\]
where \(\tau_t\) is the processing time of test \(t\). Voltage-sensitive workload \(W_e^V(S)\) is defined similarly by summing only products with \(\nu_p=1\).

This formulation captures the central trade-off. Increasing \(s_p\) increases pull-down workload but may reduce tardiness by enabling downstream parallelism. Reducing \(s_p\) saves workload but may serialize downstream testing on too few physical units.

\subsection{Feasibility conditions and objective}
\label{subsec:objective}

A feasible schedule must satisfy the following conditions:
\begin{enumerate}
    \item all and only required tests are scheduled with the multiplicities induced by \(S\);
    \item each operation is assigned to a chamber whose environment matches the test requirement;
    \item humidity-controlled tests are assigned only to chambers with \(h_c=1\);
    \item if \(\nu_p=1\), operations of product \(p\) are assigned only to chambers with \(v_c=1\);
    \item each station processes at most one operation at a time;
    \item each physical sample is used by at most one operation at a time;
    \item stage-state logic is satisfied: gas precedes pull-down, downstream tests are released after pull-down, and consumer-usage tests require the relevant downstream-stage conditions;
    \item product completion time is the finish time of its last required operation.
\end{enumerate}

Let \(C_p\) be the completion time of product \(p\). Product tardiness is
\[
T_p=\max(0,C_p-d_p),
\]
and the primary objective is
\[
\min T^{tot}=\sum_{p\in\mathcal{P}}T_p.
\]
Among solutions with equal total tardiness, the solution with smaller total sample usage
\[
S^{tot}=\sum_{p\in\mathcal{P}}s_p
\]
is preferred. The evaluation tuple is therefore \(Eval=(T^{tot},S^{tot})\) under lexicographic ordering.

\begin{table}[H]
\centering
\caption{Main notation.}
\label{tab:notation_short}
\small
\begin{tabular}{ll}
\toprule
\textbf{Symbol} & \textbf{Definition} \\
\midrule
\(\mathcal{P},\mathcal{C},\mathcal{T},\mathcal{E}\) & Products, chambers, tests, and environments \\
\(M_c,h_c,v_c\) & Stations, humidity capability, and voltage capability of chamber \(c\) \\
\(\mathcal{T}_p,d_p,\nu_p\) & Required tests, due date, and voltage requirement of product \(p\) \\
\(s_p,S,s_{\min},s_{\max}\) & Sample count, sample vector, and sample bounds \\
\(e(t),\tau_t,m_{pt}(s_p)\) & Test environment, duration, and induced multiplicity \\
\(X_c,X\) & Environment assigned to chamber \(c\), and assignment vector \\
\(W_e(S),W_e^V(S)\) & Total and voltage-sensitive workload in environment \(e\) \\
\(C_p,T_p,T^{tot},S^{tot}\) & Completion time, tardiness, total tardiness, and total samples \\
\bottomrule
\end{tabular}
\end{table}

\section{Solution Methodology}
\label{sec:solution}

\subsection{Overview}
\label{subsec:framework}

The proposed solution framework coordinates three decision layers: chamber-environment assignment, sample-size optimization, and feasible scheduling. The method starts from the fixed three-sample policy. At each outer iteration, chamber settings are selected for the current sample vector, sample counts are improved under the fixed chamber assignment, and the resulting solution is evaluated by a constructive scheduler. If local improvement stagnates, a VNS diversification phase perturbs sample sizes and re-applies local improvement.

\begin{algorithm}[H]
\caption{Integrated TCSSOP heuristic}
\label{alg:integrated_short}
\begin{algorithmic}[1]
\State Initialize \(S^{current}\) with \(s_p=3\) for all products.
\State Set incumbent \((X^{best},S^{best},\Pi^{best},Eval^{best})\leftarrow(\emptyset,\emptyset,\emptyset,(+\infty,+\infty))\).
\For{\(iter=1,\ldots,I_{\max}\)}
    \State Assign chamber environments \(X^{current}\) for \(S^{current}\).
    \State Apply sample-size local search under \(X^{current}\), obtaining \((S^{local},\Pi^{local},Eval^{local})\).
    \If{\(Eval^{local}\) improves \(Eval^{best}\)}
        \State Update incumbent and continue with \(S^{current}\leftarrow S^{local}\).
    \Else
        \State Apply VNS diversification to \(S^{best}\) under \(X^{best}\).
        \If{VNS improves the incumbent}
            \State Update incumbent and continue.
        \Else
            \State Terminate.
        \EndIf
    \EndIf
\EndFor
\State \Return incumbent solution.
\end{algorithmic}
\end{algorithm}

\subsection{Chamber assignment}
\label{subsec:chamber_assignment}

For a fixed sample vector \(S\), chamber environments are determined from environment-specific workload shares. Let
\[
W^{tot}(S)=\sum_{e\in\mathcal{E}}W_e(S),
\qquad
W^{V,tot}(S)=\sum_{e\in\mathcal{E}}W_e^V(S).
\]
The workload shares are
\[
\phi_e=\frac{W_e(S)}{W^{tot}(S)},\qquad
\phi_e^V=\frac{W_e^V(S)}{W^{V,tot}(S)}.
\]
With \(M^{tot}=\sum_{c\in\mathcal{C}}M_c\), target capacities are \(T_e=\phi_e M^{tot}\) and \(T_e^V=\phi_e^V M^{tot}\).

Let \(x_{ce}=1\) if chamber \(c\) is assigned to environment \(e\), and let \(\mathcal{E}^H\) be the set of humidity-controlled environments. The exact assignment model minimizes absolute deviations from workload-based target capacities:
\begin{align}
\min\quad & \sum_{e\in\mathcal{E}}z_e+\sum_{e\in\mathcal{E}}z_e^V \\
\text{s.t.}\quad
& \sum_{e\in\mathcal{E}}x_{ce}=1 && \forall c\in\mathcal{C},\\
& \sum_{c\in\mathcal{C}}x_{ce}\ge 1 && \forall e\in\mathcal{E},\\
& x_{ce}\le h_c && \forall c\in\mathcal{C},\ e\in\mathcal{E}^H,\\
& z_e\ge \sum_{c\in\mathcal{C}}M_cx_{ce}-T_e && \forall e\in\mathcal{E},\\
& z_e\ge T_e-\sum_{c\in\mathcal{C}}M_cx_{ce} && \forall e\in\mathcal{E},\\
& z_e^V\ge \sum_{c\in\mathcal{C}}M_cv_cx_{ce}-T_e^V && \forall e\in\mathcal{E},\\
& z_e^V\ge T_e^V-\sum_{c\in\mathcal{C}}M_cv_cx_{ce} && \forall e\in\mathcal{E},\\
& x_{ce}\in\{0,1\},\quad z_e,z_e^V\ge 0 && \forall c,e.
\end{align}
The resulting assignment is then refined by swap and move neighborhoods. A candidate is retained only if it preserves environment coverage and capability feasibility. Feasible candidates are evaluated with a fast schedule-aware scoring function, and a first-improvement rule is used.

\subsection{Sample-size local search}
\label{subsec:sample_optimization}

Given chamber assignment \(X\), local search explores product-level sample changes. For each product \(p\), the neighborhood is
\[
\Delta_p \subseteq \{-2,-1,+1,+2\}
\]
subject to sample bounds. Candidate vectors are first screened by fast schedule-aware evaluation. If a move improves the lexicographic objective, the full scheduler is executed to obtain the exact schedule and objective value.

A Late--Early Rebalance operator is also used. Products with positive tardiness are paired with products having earliness slack. For each pair, one sample is added to the late product and one sample is removed from the early product, subject to bounds. The paired move is accepted if total tardiness does not worsen and, under equal tardiness, total sample usage does not increase.

\subsection{Scheduling and feasibility validation}
\label{subsec:scheduling}

For fixed \(X\) and \(S\), schedules are constructed using an earliest-due-date (EDD) ready-set rule. Candidate jobs are generated project-wise from the current stage state, sample availability, and earliest feasible station time. The best ready job of each product enters a global candidate pool, and the next job is selected by due date, with earliest feasible start time as a tie-breaker.

Each selected job is assigned to the earliest compatible chamber station. Compatibility requires matching environment, humidity capability when needed, and voltage capability for voltage-sensitive products. Sample availability is updated after every assignment to prevent overlap on the same physical unit.

After construction, the schedule is validated independently. The validation checks required-test multiplicity, chamber-environment compatibility, humidity and voltage feasibility, station non-overlap, sample non-overlap, stage-state logic, and completion-time calculation. This layer is essential because infeasible schedules are operationally unusable regardless of objective value.

\subsection{VNS diversification}
\label{subsec:vns}

VNS is activated when the outer loop does not improve the incumbent. The method uses a fixed fallback ratio sequence
\[
R^{VNS}=[0.20,0.30].
\]
For each ratio, up to \(K_{\max}\) shake attempts are performed. A shake selects the corresponding fraction of products and alternates between two directions: \textsc{DOWN}, which resets selected samples to \(s_{\min}\), and \textsc{UP}, which resets selected samples to the VNS upper level. The perturbed vector is then improved by the same sample-size local search. The first locally improved candidate that improves the incumbent is accepted.

This design changes the workload structure rather than only the sequence of existing operations. As a result, VNS can escape local optima created by the current balance between pull-down workload, downstream parallelism, and chamber eligibility.

\section{Computational Study}
\label{sec:computational_study}

\subsection{Experimental design}
\label{subsec:exp_design}

The benchmark set contains 90 scenario combinations formed by \(10\) test matrices, \(3\) due-date scenarios, and \(3\) voltage scenarios. The instances are synthetically generated using ranges and structural rules calibrated from industrial practice, allowing controlled experiments without relying on proprietary raw data.

Due-date scenarios represent tight, moderate, and loose planning regimes with different dispersion levels. Voltage scenarios define low, medium, and high proportions of voltage-sensitive products; higher voltage intensity reduces the eligible chamber set and increases competition for voltage-capable resources. Test matrices represent different product portfolios and required-test combinations. They mainly affect workload distribution across environments because gas and pull-down tests concentrate upstream demand, while energy, performance, freezing, temperature-rise, and consumer-usage tests distribute downstream workload across several temperatures.

\begin{figure}[H]
\centering
\includegraphics[width=0.72\textwidth]{duedate_senaryo.png}
\caption{Due-date distributions across scenarios.}
\label{fig:due_date_spread_short}
\end{figure}

\begin{figure}[H]
\centering
\includegraphics[width=0.66\textwidth]{voltage_senaryo.png}
\caption{Voltage requirement intensity across scenarios.}
\label{fig:voltage_intensity_short}
\end{figure}

\subsection{Compared methods and parameters}
\label{subsec:methods_params}

Four variants are compared:
\begin{itemize}
    \item \textbf{FS (Fixed Sample):} all products use \(s_p=3\); chamber assignment is performed once and no sample optimization is applied.
    \item \textbf{SO-noOuter:} sample sizes are optimized under the initial chamber assignment, without iterative chamber reassignment.
    \item \textbf{SO-noVNS:} sample optimization and chamber reassignment are iterated, but VNS diversification is disabled.
    \item \textbf{Proposed Framework:} sample optimization, chamber reassignment, and VNS diversification are all enabled.
\end{itemize}

All parameters are fixed across instances. The initial sample count is \(s_p^{(0)}=3\), \(s_{\min}=3\), \(s_{\max}=11\), and the VNS \textsc{UP} level is \(s_{\max VNS}=8\). The sample neighborhood is \(\{-2,-1,+1,+2\}\), the VNS fallback ratios are \([0.20,0.30]\), \(I_{\max}=100\), and \(K_{\max}=200\) for the selected configuration. The objective is evaluated lexicographically as \((T^{tot},S^{tot})\). Runtime is reported separately.

\subsection{Performance measures}
\label{subsec:metrics}

For method \(m\) and instance \(i\), the percentage gap is computed from total tardiness:
\[
Gap_{im}=
\frac{T^{tot}_{im}-T^{tot,best}_{i}}
{T^{tot,best}_{i}}\times 100,
\]
where \(T^{tot,best}_{i}\) is the best total tardiness observed for that instance in the relevant comparison. In the reported experiments, all denominators are positive. Best-instance counts are based on the lexicographic objective, so total sample usage breaks ties when total tardiness is equal.

\subsection{Results}
\label{subsec:results}

Table~\ref{tab:vns_ratio_comparison_short} compares VNS perturbation policies with \(K_{\max}=100\). The \(0.20\)--\(0.30\) fallback policy produces the lowest average gap and standard deviation among these variants, showing that moderate perturbation followed by stronger fallback is more robust than a single fixed ratio.

\begin{table}[H]
\centering
\caption{Comparison of VNS perturbation policies with \(K_{\max}=100\).}
\label{tab:vns_ratio_comparison_short}
\small
\begin{tabular}{lccccc}
\toprule
\textbf{Measure} &
\makecell{\textbf{0.10}\\\textbf{VNS}} &
\makecell{\textbf{0.20}\\\textbf{VNS}} &
\makecell{\textbf{0.30}\\\textbf{VNS}} &
\makecell{\textbf{0.10--0.20--}\\\textbf{0.30 VNS}} &
\makecell{\textbf{0.20--0.30}\\\textbf{VNS}} \\
\midrule
Average gap (\%) & 29.62 & 18.11 & 23.44 & 22.06 & 14.93 \\
Std. dev. (\%) & 44.09 & 31.52 & 30.10 & 33.25 & 27.35 \\
Best-instance count & 25 & 27 & 18 & 38 & 36 \\
Average runtime (min) & 41.3 & 32.1 & 30.4 & 66.0 & 49.7 \\
\bottomrule
\end{tabular}
\end{table}

Table~\ref{tab:vns_budget_comparison_short} evaluates the shake budget for the selected \(0.20\)--\(0.30\) policy. Increasing \(K_{\max}\) improves both average performance and robustness. With \(K_{\max}=200\), the average gap is only \(0.31\%\), the standard deviation is \(1.65\%\), and the method obtains the best result in 82 of 90 scenarios within this comparison. The improvement comes at higher runtime, which is acceptable for tactical planning.

\begin{table}[H]
\centering
\caption{Effect of shake budget for the \(0.20\)--\(0.30\) VNS policy.}
\label{tab:vns_budget_comparison_short}
\begin{tabular}{lccc}
\toprule
\textbf{Measure} &
\makecell{\textbf{VNS}\\\textbf{(50)}} &
\makecell{\textbf{VNS}\\\textbf{(100)}} &
\makecell{\textbf{VNS}\\\textbf{(200)}} \\
\midrule
Average gap (\%) & 13.79 & 5.65 & 0.31 \\
Std. dev. (\%) & 29.07 & 20.06 & 1.65 \\
Best-instance count & 18 & 47 & 82 \\
Average runtime (min) & 32.6 & 49.7 & 85.1 \\
\bottomrule
\end{tabular}
\end{table}

The selected proposed configuration is then compared with the main benchmark variants in Table~\ref{tab:main_benchmark_comparison_short}. It obtains the best solution in all 90 scenarios. The fixed-sample policy has the largest average gap, confirming that a uniform sample rule is not robust under heterogeneous due dates, environmental requirements, and chamber-capability constraints. The no-outer and no-VNS variants also remain far from the selected configuration, showing that both assignment--scheduling feedback and diversification are necessary.

\begin{table}[H]
\centering
\caption{Overall comparison of the selected proposed configuration and benchmark variants.}
\label{tab:main_benchmark_comparison_short}
\begin{tabular}{lcccc}
\toprule
\textbf{Method} &
\makecell{\textbf{Average}\\\textbf{gap (\%)}} &
\makecell{\textbf{Std. dev.}\\\textbf{(\%)}} &
\makecell{\textbf{Best-instance}\\\textbf{count}} &
\makecell{\textbf{Average runtime}\\\textbf{(min)}} \\
\midrule
Fixed sample & 361.93 & 198.43 & 0 & 0.1 \\
No outer reconfiguration & 140.07 & 90.32 & 0 & 17.1 \\
No VNS diversification & 178.67 & 113.52 & 0 & 0.5 \\
Selected proposed configuration & 0.00 & 0.00 & 90 & 85.1 \\
\bottomrule
\end{tabular}
\end{table}

Table~\ref{tab:fixed_selected_sample_comparison_short} provides a direct comparison between the fixed policy and the selected configuration. The selected configuration uses more samples on average, but it allocates them selectively to products for which additional parallelism reduces tardiness. The goal is therefore not uniform prototype expansion, but coordinated sample allocation and chamber reconfiguration.

\begin{table}[H]
\centering
\caption{Direct comparison between fixed samples and the selected proposed configuration.}
\label{tab:fixed_selected_sample_comparison_short}
\begin{tabular}{lcc}
\toprule
\textbf{Measure} & \textbf{Fixed sample} & \makecell{\textbf{Selected proposed}\\\textbf{configuration}} \\
\midrule
Average gap (\%) & 361.93 & 0.00 \\
Std. dev. (\%) & 198.43 & 0.00 \\
Best-instance count & 0 & 90 \\
Average runtime (min) & 0.1 & 85.1 \\
Average sample count & 150 & 380 \\
\bottomrule
\end{tabular}
\end{table}

\subsection{Managerial implications}
\label{subsec:managerial}

The results provide three managerial insights. First, sample counts should be optimized rather than fixed. The value of an additional sample depends on due-date urgency, downstream parallelization potential, and the compatible chamber set. Second, sample optimization and chamber assignment should be updated together because changing sample counts shifts environment-specific and voltage-sensitive workload. Third, diversification is important in tactical planning: local sample changes alone may not escape a poor workload structure, whereas controlled VNS perturbations can reveal substantially better sample allocations.

\section{Conclusion}
\label{sec:conc}

This study introduced TCSSOP, an integrated test-chamber scheduling problem with decision-dependent workload generation. The problem arises in refrigerator R\&D testing, where heterogeneous chambers, environmental compatibility, voltage and humidity restrictions, stage-state logic, due dates, and physical sample availability interact.

The proposed framework jointly optimizes chamber settings, sample sizes, and feasible schedules. Chamber environments are assigned using workload-balancing logic; sample sizes are improved by local search and Late--Early Rebalance; schedules are constructed by an EDD-based feasible scheduler; and VNS diversification perturbs sample decisions when local improvement stagnates. Feasibility is enforced through explicit checks for required tests, chamber compatibility, humidity and voltage eligibility, station and sample non-overlap, and stage-state correctness.

Computational experiments on 90 benchmark scenarios show that the selected proposed configuration with \(0.20\)--\(0.30\) VNS fallback and \(K_{\max}=200\) consistently outperforms fixed-sample and ablated variants. The fixed three-sample policy is fast but produces substantially higher tardiness. The ablation results show that both iterative chamber reconfiguration and VNS diversification materially improve solution quality.

The main practical conclusion is that prototype sample allocation should be treated as a tactical planning variable. Samples should be increased selectively where additional parallelism reduces tardiness, and chamber configurations should be reassessed after sample decisions change. Future work may incorporate uncertain test durations, test failures and retesting, chamber setup times, energy and labor objectives, and rolling-horizon project arrivals.

\bibliographystyle{plainnat}
\bibliography{References}

\end{document}
